%% Generated by Sphinx.
\def\sphinxdocclass{report}
\documentclass[letterpaper,10pt,english]{sphinxmanual}
\ifdefined\pdfpxdimen
   \let\sphinxpxdimen\pdfpxdimen\else\newdimen\sphinxpxdimen
\fi \sphinxpxdimen=.75bp\relax

\PassOptionsToPackage{warn}{textcomp}
\usepackage[utf8]{inputenc}
\ifdefined\DeclareUnicodeCharacter
% support both utf8 and utf8x syntaxes
\edef\sphinxdqmaybe{\ifdefined\DeclareUnicodeCharacterAsOptional\string"\fi}
  \DeclareUnicodeCharacter{\sphinxdqmaybe00A0}{\nobreakspace}
  \DeclareUnicodeCharacter{\sphinxdqmaybe2500}{\sphinxunichar{2500}}
  \DeclareUnicodeCharacter{\sphinxdqmaybe2502}{\sphinxunichar{2502}}
  \DeclareUnicodeCharacter{\sphinxdqmaybe2514}{\sphinxunichar{2514}}
  \DeclareUnicodeCharacter{\sphinxdqmaybe251C}{\sphinxunichar{251C}}
  \DeclareUnicodeCharacter{\sphinxdqmaybe2572}{\textbackslash}
\fi
\usepackage{cmap}
\usepackage[T1]{fontenc}
\usepackage{amsmath,amssymb,amstext}
\usepackage{babel}
\usepackage{times}
\usepackage[Sonny]{fncychap}
\ChNameVar{\Large\normalfont\sffamily}
\ChTitleVar{\Large\normalfont\sffamily}
\usepackage{sphinx}

\fvset{fontsize=\small}
\usepackage{geometry}

% Include hyperref last.
\usepackage{hyperref}
% Fix anchor placement for figures with captions.
\usepackage{hypcap}% it must be loaded after hyperref.
% Set up styles of URL: it should be placed after hyperref.
\urlstyle{same}

\addto\captionsenglish{\renewcommand{\figurename}{Figura }}
\makeatletter
\def\fnum@figure{\figurename\thefigure{}}
\makeatother
\addto\captionsenglish{\renewcommand{\tablename}{Tabla }}
\makeatletter
\def\fnum@table{\tablename\thetable{}}
\makeatother
\addto\captionsenglish{\renewcommand{\literalblockname}{Lista}}

\addto\captionsenglish{\renewcommand{\literalblockcontinuedname}{proviene de la página anterior}}
\addto\captionsenglish{\renewcommand{\literalblockcontinuesname}{continué en la próxima página}}
\addto\captionsenglish{\renewcommand{\sphinxnonalphabeticalgroupname}{Non-alphabetical}}
\addto\captionsenglish{\renewcommand{\sphinxsymbolsname}{Símbolos}}
\addto\captionsenglish{\renewcommand{\sphinxnumbersname}{Numbers}}

\addto\extrasenglish{\def\pageautorefname{página}}

\setcounter{tocdepth}{1}



\title{Escaleras y Serpientes Documentation}
\date{08 de July de 2020}
\release{1}
\author{Daniel Salazar, Josue Amador, Gustavo Jimenez, Ricardo Franco, Ignacio Arroyo}
\newcommand{\sphinxlogo}{\vbox{}}
\renewcommand{\releasename}{Versión}
\makeindex
\begin{document}

\ifdefined\shorthandoff
  \ifnum\catcode`\=\string=\active\shorthandoff{=}\fi
  \ifnum\catcode`\"=\active\shorthandoff{"}\fi
\fi

\pagestyle{empty}
\sphinxmaketitle
\pagestyle{plain}
\sphinxtableofcontents
\pagestyle{normal}
\phantomsection\label{\detokenize{Escaleras_y_Serpientes::doc}}



\chapter{\sphinxstylestrong{Informacion importante}}
\label{\detokenize{Escaleras_y_Serpientes:informacion-importante}}

\section{\sphinxstylestrong{Codigo Generado automaticamente}}
\label{\detokenize{codigo:codigo-generado-automaticamente}}\label{\detokenize{codigo::doc}}

\subsection{Codigo de: \sphinxstylestrong{Jugador}}
\label{\detokenize{codigo:module-jugador}}\label{\detokenize{codigo:codigo-de-jugador}}\index{jugador (módulo)@\spxentry{jugador}\spxextra{módulo}}\index{Jugador (clase en jugador)@\spxentry{Jugador}\spxextra{clase en jugador}}

\begin{fulllineitems}
\phantomsection\label{\detokenize{codigo:jugador.Jugador}}\pysiglinewithargsret{\sphinxbfcode{\sphinxupquote{class }}\sphinxcode{\sphinxupquote{jugador.}}\sphinxbfcode{\sphinxupquote{Jugador}}}{\emph{pos}, \emph{x}, \emph{y}, \emph{a}, \emph{l}, \emph{mX}, \emph{mY}}{}
Esta clase corresponde a la del Jugador que participa en el juego. Realiza las 
interacciones dadas por el usuario y depende totalmente de el. Van a haber entre 2 
a 4 jugadores por juego.
\index{getNombre() (método de jugador.Jugador)@\spxentry{getNombre()}\spxextra{método de jugador.Jugador}}

\begin{fulllineitems}
\phantomsection\label{\detokenize{codigo:jugador.Jugador.getNombre}}\pysiglinewithargsret{\sphinxbfcode{\sphinxupquote{getNombre}}}{}{}
Metodo que va a devolver el nombre del Jugador.
:returns: nombre del jugador, es tipo string.

\end{fulllineitems}

\index{mover() (método de jugador.Jugador)@\spxentry{mover()}\spxextra{método de jugador.Jugador}}

\begin{fulllineitems}
\phantomsection\label{\detokenize{codigo:jugador.Jugador.mover}}\pysiglinewithargsret{\sphinxbfcode{\sphinxupquote{mover}}}{\emph{distancia}, \emph{coordenadas}}{}
Metodo que se encarga de realizar el movimiento del jugador y requiere de la participacion del
Arbitro para saber si: gano; si cayo en una escalera o serpiente y necesita moverse (tampoco sabe
adonde habria que moverse, el Arbitro le indica);y por ultimo, si ya no tiene que moverse mas.
:param distancia: cantidad de casillas a moverse, puede ser el numero que salio en los dados, o
puede ser lo que debe moverse por haber caido en escalera o serpiente.
:type distancia: int
:param coordenadas: coordenadas en X y Y de la casilla (especificamente en el centro de ella) adonde 
el Jugador tiene que moverse.
:type coordenadas: tupla de int

\end{fulllineitems}

\index{setNombre() (método de jugador.Jugador)@\spxentry{setNombre()}\spxextra{método de jugador.Jugador}}

\begin{fulllineitems}
\phantomsection\label{\detokenize{codigo:jugador.Jugador.setNombre}}\pysiglinewithargsret{\sphinxbfcode{\sphinxupquote{setNombre}}}{\emph{nombre}}{}
Metodo que se encarga de darle nombre al Jugador y almacenarlo en la variable nombre.
:param nombre: nombre para ponerle al Jugador
:type nombre: string

\end{fulllineitems}


\end{fulllineitems}

\index{VistaJugador (clase en jugador)@\spxentry{VistaJugador}\spxextra{clase en jugador}}

\begin{fulllineitems}
\phantomsection\label{\detokenize{codigo:jugador.VistaJugador}}\pysiglinewithargsret{\sphinxbfcode{\sphinxupquote{class }}\sphinxcode{\sphinxupquote{jugador.}}\sphinxbfcode{\sphinxupquote{VistaJugador}}}{\emph{jugador}}{}
Esta clase corresponde a la Vista del Jugador que participa en el juego. Lo va a dibujar en la pantalla de pygame
para que pueda ser visible para los usuarios.
\index{dibujar() (método de jugador.VistaJugador)@\spxentry{dibujar()}\spxextra{método de jugador.VistaJugador}}

\begin{fulllineitems}
\phantomsection\label{\detokenize{codigo:jugador.VistaJugador.dibujar}}\pysiglinewithargsret{\sphinxbfcode{\sphinxupquote{dibujar}}}{}{}
Metodo que ilustra al jugador en la vetana de pygame. Lo dibuja para que el usuario
pueda observarlo.

\end{fulllineitems}

\index{muestre\_turno\_jugador() (método de jugador.VistaJugador)@\spxentry{muestre\_turno\_jugador()}\spxextra{método de jugador.VistaJugador}}

\begin{fulllineitems}
\phantomsection\label{\detokenize{codigo:jugador.VistaJugador.muestre_turno_jugador}}\pysiglinewithargsret{\sphinxbfcode{\sphinxupquote{muestre\_turno\_jugador}}}{}{}
Este metodo se va a encargar de desplegar en la ventana de pygame que jugador
le toca jugar. El texto lo muestra en el color respectivo de cada jugador. 
:param jugador: Objeto de jugador a realizar turno.
:type jugador: Jugador object

\end{fulllineitems}


\end{fulllineitems}



\subsection{Codigo de: \sphinxstylestrong{Arbitro}}
\label{\detokenize{codigo:module-arbitro}}\label{\detokenize{codigo:codigo-de-arbitro}}\index{arbitro (módulo)@\spxentry{arbitro}\spxextra{módulo}}\index{Arbitro (clase en arbitro)@\spxentry{Arbitro}\spxextra{clase en arbitro}}

\begin{fulllineitems}
\phantomsection\label{\detokenize{codigo:arbitro.Arbitro}}\pysigline{\sphinxbfcode{\sphinxupquote{class }}\sphinxcode{\sphinxupquote{arbitro.}}\sphinxbfcode{\sphinxupquote{Arbitro}}}
Clase de Arbitro que se encarga de revisar el comportamiento de los jugadores. Se encarga de
asegurar el buen funcionamiento del juego y lo realiza de las siguientes maneras: comprueba si 
un jugador gano; si se tiene que mover porque cayo en escaelera o serpiente (si lo hace, le dice
hacia donde moverse), ayuda a decidir los turnos al inicio del juego, reporta los jugadores como
existentes dentro del juego, se asegura que cada jugador tenga un color unico...
\index{aplicarReglasMovimiento() (método de clase de arbitro.Arbitro)@\spxentry{aplicarReglasMovimiento()}\spxextra{método de clase de arbitro.Arbitro}}

\begin{fulllineitems}
\phantomsection\label{\detokenize{codigo:arbitro.Arbitro.aplicarReglasMovimiento}}\pysiglinewithargsret{\sphinxbfcode{\sphinxupquote{classmethod }}\sphinxbfcode{\sphinxupquote{aplicarReglasMovimiento}}}{\emph{posTablero}}{}
Este metodo comprueba la posicion del jugador en el tablero y si en esa casilla hay una serpiente o escalera, y aplica el movimiento
indicado
:param posTablero: recibe la posicion del jugador en el tablero
:type posTablero: numero entero
:return distancia: retorna la distancia que se tiene que recorrer, ya sea al final de la serpiente o el tope de la escalera
:return coordenadasNuevas: retorna las coordenadas de la nueva casilla hacia donde se va a mover el jugador

\end{fulllineitems}

\index{buscar\_jugador() (método de clase de arbitro.Arbitro)@\spxentry{buscar\_jugador()}\spxextra{método de clase de arbitro.Arbitro}}

\begin{fulllineitems}
\phantomsection\label{\detokenize{codigo:arbitro.Arbitro.buscar_jugador}}\pysiglinewithargsret{\sphinxbfcode{\sphinxupquote{classmethod }}\sphinxbfcode{\sphinxupquote{buscar\_jugador}}}{\emph{jugadores}, \emph{valor}}{}
Este metodo busca un jugador dentro del vector de turnos de acuerdo al valor que saco en los dados
:param jugadores: el vector que guarda los turnos de los jugadores
:type jugadores: un array de Jugador
:param valor: la cantidad que tiro en los dados el jugador
:type valor: numero entero

\end{fulllineitems}

\index{comprobarGane() (método de clase de arbitro.Arbitro)@\spxentry{comprobarGane()}\spxextra{método de clase de arbitro.Arbitro}}

\begin{fulllineitems}
\phantomsection\label{\detokenize{codigo:arbitro.Arbitro.comprobarGane}}\pysiglinewithargsret{\sphinxbfcode{\sphinxupquote{classmethod }}\sphinxbfcode{\sphinxupquote{comprobarGane}}}{\emph{pos}}{}
Este metodo se encarga de verificar si un jugador ha ganado la partida, se utiliza en cada jugada del jugador
:param pos: Recibe la posicion de la casilla en la que se encuentra el jugador 
:type pos: numero entero
:returns: bool, devuelve true si el jugador en cuestion gano la partida

\end{fulllineitems}

\index{decidir\_turnos() (método de clase de arbitro.Arbitro)@\spxentry{decidir\_turnos()}\spxextra{método de clase de arbitro.Arbitro}}

\begin{fulllineitems}
\phantomsection\label{\detokenize{codigo:arbitro.Arbitro.decidir_turnos}}\pysiglinewithargsret{\sphinxbfcode{\sphinxupquote{classmethod }}\sphinxbfcode{\sphinxupquote{decidir\_turnos}}}{\emph{cantidad\_jugadores}}{}
Este metodo se utiliza para decidir el orden de los turnos de los jugadores, cada uno lanza dados y el mas alto, va de primero y asi sucesivamente
:param cantidad\_jugadores: recibe la cantidad de jugadores que estan en la partida
:type cantidad\_jugadores: numero entero
:param boton\_dados: recibe una instancia de un boton de dados del cual se sacara la informacion de los dados
:type boton\_dados: instancia de la clase button
:param dado: el dado que se lanzara para sacar un numero al azar
:type dado: una instancia de dado

\end{fulllineitems}

\index{esColorRepetido() (método de clase de arbitro.Arbitro)@\spxentry{esColorRepetido()}\spxextra{método de clase de arbitro.Arbitro}}

\begin{fulllineitems}
\phantomsection\label{\detokenize{codigo:arbitro.Arbitro.esColorRepetido}}\pysiglinewithargsret{\sphinxbfcode{\sphinxupquote{classmethod }}\sphinxbfcode{\sphinxupquote{esColorRepetido}}}{\emph{color}}{}
Va a buscar si el color, ya acordado por el usuario, le pertenece a otro jugador y retorna true si es asi.
Si por el otro lado, no esta repetido, entonces retorna false.
:param color: color a buscar dentro de los colores de jugadores
:type color: string
:returns: True si un jugador ya tiene ese color, False si nadie lo ha solicitado antes y no le pertenece a alguien.

\end{fulllineitems}

\index{existeColor() (método de clase de arbitro.Arbitro)@\spxentry{existeColor()}\spxextra{método de clase de arbitro.Arbitro}}

\begin{fulllineitems}
\phantomsection\label{\detokenize{codigo:arbitro.Arbitro.existeColor}}\pysiglinewithargsret{\sphinxbfcode{\sphinxupquote{classmethod }}\sphinxbfcode{\sphinxupquote{existeColor}}}{\emph{color}}{}
Va a buscar dentro de los colores existentes, si el color escogido por el usuario existe, y va a 
retornar True si si exsite y False si no.
:param color: Color escogido por el usuario que se va a buscar si existe
:type color: string
:returns: retorna True si encuentra el color dentro de los existentes, False si no.

\end{fulllineitems}

\index{init() (método de clase de arbitro.Arbitro)@\spxentry{init()}\spxextra{método de clase de arbitro.Arbitro}}

\begin{fulllineitems}
\phantomsection\label{\detokenize{codigo:arbitro.Arbitro.init}}\pysiglinewithargsret{\sphinxbfcode{\sphinxupquote{classmethod }}\sphinxbfcode{\sphinxupquote{init}}}{\emph{controlador}}{}
Inicializa el Arbitro.
:param controlador: referencia al conrolador.

\end{fulllineitems}

\index{jugador\_juegue() (método de clase de arbitro.Arbitro)@\spxentry{jugador\_juegue()}\spxextra{método de clase de arbitro.Arbitro}}

\begin{fulllineitems}
\phantomsection\label{\detokenize{codigo:arbitro.Arbitro.jugador_juegue}}\pysiglinewithargsret{\sphinxbfcode{\sphinxupquote{classmethod }}\sphinxbfcode{\sphinxupquote{jugador\_juegue}}}{\emph{jugador}}{}
Metodo que inicia el turno de un jugador, y espera al evento que este tire los dados.

\end{fulllineitems}

\index{obtener\_destino() (método de clase de arbitro.Arbitro)@\spxentry{obtener\_destino()}\spxextra{método de clase de arbitro.Arbitro}}

\begin{fulllineitems}
\phantomsection\label{\detokenize{codigo:arbitro.Arbitro.obtener_destino}}\pysiglinewithargsret{\sphinxbfcode{\sphinxupquote{classmethod }}\sphinxbfcode{\sphinxupquote{obtener\_destino}}}{\emph{index}, \emph{resultado\_dados}}{}
Metodo que va a conseguir las coordenadas de la casilla adonde el jugador tiene
que moverse. 
:param index: posicion en el tablero donde esta el usuario
:type index: int
:param resultado\_dados: cantidad de casillas a moverse
:type resultado\_dados: int

\end{fulllineitems}

\index{procesaRespuestaDados() (método de clase de arbitro.Arbitro)@\spxentry{procesaRespuestaDados()}\spxextra{método de clase de arbitro.Arbitro}}

\begin{fulllineitems}
\phantomsection\label{\detokenize{codigo:arbitro.Arbitro.procesaRespuestaDados}}\pysiglinewithargsret{\sphinxbfcode{\sphinxupquote{classmethod }}\sphinxbfcode{\sphinxupquote{procesaRespuestaDados}}}{\emph{resultado\_dados}}{}
Este metodo se encarga de leer el resultado de dados, para decidir el movimiento del jugador.
Tambien indica el final de la jugada del jugador, y el Arbitro analizara posteriormente el
movimiento del jugador.
:param resultado\_dados: resultado de dados, se almacenara en self.valor\_actual\_dados.

\end{fulllineitems}

\index{reciboRespuestaDados() (método de clase de arbitro.Arbitro)@\spxentry{reciboRespuestaDados()}\spxextra{método de clase de arbitro.Arbitro}}

\begin{fulllineitems}
\phantomsection\label{\detokenize{codigo:arbitro.Arbitro.reciboRespuestaDados}}\pysiglinewithargsret{\sphinxbfcode{\sphinxupquote{classmethod }}\sphinxbfcode{\sphinxupquote{reciboRespuestaDados}}}{\emph{resultado\_dados}}{}
Este metodo se encarga de leer el resultado de dados, para decidir el orden de turnos 
de los jugadores.
:param resultado\_dados: resultado de dados, se almacenara en un vector el jugador y su
resultado especifico.

\end{fulllineitems}

\index{reporte\_jugadores() (método de clase de arbitro.Arbitro)@\spxentry{reporte\_jugadores()}\spxextra{método de clase de arbitro.Arbitro}}

\begin{fulllineitems}
\phantomsection\label{\detokenize{codigo:arbitro.Arbitro.reporte_jugadores}}\pysiglinewithargsret{\sphinxbfcode{\sphinxupquote{classmethod }}\sphinxbfcode{\sphinxupquote{reporte\_jugadores}}}{\emph{jugador}}{}
Este metodo agrega un Jugador a la lista de jugadores que tiene el Singleton, la cual es informacion importante para las demas clases
:param jugador: recibe una instancia de un Jugador para agregarlo a la lista
:type jugador: instancia de la clase jugador

\end{fulllineitems}

\index{setColorJugador() (método de clase de arbitro.Arbitro)@\spxentry{setColorJugador()}\spxextra{método de clase de arbitro.Arbitro}}

\begin{fulllineitems}
\phantomsection\label{\detokenize{codigo:arbitro.Arbitro.setColorJugador}}\pysiglinewithargsret{\sphinxbfcode{\sphinxupquote{classmethod }}\sphinxbfcode{\sphinxupquote{setColorJugador}}}{\emph{jugador}}{}
Va a almacenar un color para el Jugador y esto es decision del usuario. El color tiene que existir dentro de las opciones
y no puede repetirse. El color es unico por jugador en la partida.

\end{fulllineitems}


\end{fulllineitems}



\subsection{Codigo de: \sphinxstylestrong{Dado}}
\label{\detokenize{codigo:module-dado}}\label{\detokenize{codigo:codigo-de-dado}}\index{dado (módulo)@\spxentry{dado}\spxextra{módulo}}\index{Dado (clase en dado)@\spxentry{Dado}\spxextra{clase en dado}}

\begin{fulllineitems}
\phantomsection\label{\detokenize{codigo:dado.Dado}}\pysigline{\sphinxbfcode{\sphinxupquote{class }}\sphinxcode{\sphinxupquote{dado.}}\sphinxbfcode{\sphinxupquote{Dado}}}
Esta clase representa un dado de seis caras, y sus atributos y acciones necesarios 
para su comportamiento en un tablero de juego.
\index{lanzar() (método de dado.Dado)@\spxentry{lanzar()}\spxextra{método de dado.Dado}}

\begin{fulllineitems}
\phantomsection\label{\detokenize{codigo:dado.Dado.lanzar}}\pysiglinewithargsret{\sphinxbfcode{\sphinxupquote{lanzar}}}{}{}
Esta función dibuja los dados resultantes de un lanzamiento y arroja la suma de sus valores.
:returns:  int -- Suma de las caras resultantes de ambos dados.

\end{fulllineitems}


\end{fulllineitems}

\index{VistaDado (clase en dado)@\spxentry{VistaDado}\spxextra{clase en dado}}

\begin{fulllineitems}
\phantomsection\label{\detokenize{codigo:dado.VistaDado}}\pysiglinewithargsret{\sphinxbfcode{\sphinxupquote{class }}\sphinxcode{\sphinxupquote{dado.}}\sphinxbfcode{\sphinxupquote{VistaDado}}}{\emph{controlador}}{}
Esta clase representa un dado de seis caras, y sus atributos y acciones necesarios 
para su comportamiento en un tablero de juego.
\index{dibujar() (método de dado.VistaDado)@\spxentry{dibujar()}\spxextra{método de dado.VistaDado}}

\begin{fulllineitems}
\phantomsection\label{\detokenize{codigo:dado.VistaDado.dibujar}}\pysiglinewithargsret{\sphinxbfcode{\sphinxupquote{dibujar}}}{\emph{resultados\_por\_dado}}{}
Esta función dibuja dos dados según el resultado de sus caras.
:returns:  None -- Procedimiento sin valor de retorno.

\end{fulllineitems}


\end{fulllineitems}



\subsection{Codigo de: \sphinxstylestrong{Almacenador}}
\label{\detokenize{codigo:module-almacenador}}\label{\detokenize{codigo:codigo-de-almacenador}}\index{almacenador (módulo)@\spxentry{almacenador}\spxextra{módulo}}\index{Almacenador (clase en almacenador)@\spxentry{Almacenador}\spxextra{clase en almacenador}}

\begin{fulllineitems}
\phantomsection\label{\detokenize{codigo:almacenador.Almacenador}}\pysigline{\sphinxbfcode{\sphinxupquote{class }}\sphinxcode{\sphinxupquote{almacenador.}}\sphinxbfcode{\sphinxupquote{Almacenador}}}
Esta clase corresponde a un almacenador, cuya funcionalidad es guardar en un archivo los avances del juego de manera
que el usuario pueda continuar la partida en otro momento. Esta clase respeta el patron de disenyo Constructor.
\index{guardar() (método de almacenador.Almacenador)@\spxentry{guardar()}\spxextra{método de almacenador.Almacenador}}

\begin{fulllineitems}
\phantomsection\label{\detokenize{codigo:almacenador.Almacenador.guardar}}\pysiglinewithargsret{\sphinxbfcode{\sphinxupquote{guardar}}}{}{}
Metodo abstracto del guardado de archivos. Cada clase que haga herencia, guardara en un archivo de formato distinto.
En nuestro caso, solo lo haremos con una clase que guarde en .txt.

\end{fulllineitems}


\end{fulllineitems}

\index{AlmacenadorTXT (clase en almacenador)@\spxentry{AlmacenadorTXT}\spxextra{clase en almacenador}}

\begin{fulllineitems}
\phantomsection\label{\detokenize{codigo:almacenador.AlmacenadorTXT}}\pysiglinewithargsret{\sphinxbfcode{\sphinxupquote{class }}\sphinxcode{\sphinxupquote{almacenador.}}\sphinxbfcode{\sphinxupquote{AlmacenadorTXT}}}{\emph{nombreArchivo}}{}
Esta clase hereda de Almacenador, y demuestra el correcto funcionamiento del patron Constructor. Esta clase es especifica
para crear archivos de formato .txt donde se almacenara la informacion vital del juego, para que sea posible jugarlo luego.
\index{guardar() (método de almacenador.AlmacenadorTXT)@\spxentry{guardar()}\spxextra{método de almacenador.AlmacenadorTXT}}

\begin{fulllineitems}
\phantomsection\label{\detokenize{codigo:almacenador.AlmacenadorTXT.guardar}}\pysiglinewithargsret{\sphinxbfcode{\sphinxupquote{guardar}}}{}{}
Metodo que va guaardar en un archivo de texto toda la informacion vital para guardar la partida. Esta informacion es: la
cantidad de jugadores; el nombre, posicion en el tablero, color, ubicacion en casilla en X y ubicacion en casilla en Y de 
cada jugador; cantidad de serpientes; posicion en el tablero de inicio de sepriente y la posicion de final de la misma; 
cantidad de escaleras; posicion en el tablero de inicio de escalera y la posicion de final de la misma.

\end{fulllineitems}


\end{fulllineitems}



\subsection{Codigo de: \sphinxstylestrong{Button}}
\label{\detokenize{codigo:module-button}}\label{\detokenize{codigo:codigo-de-button}}\index{button (módulo)@\spxentry{button}\spxextra{módulo}}

\subsection{Codigo de: \sphinxstylestrong{Cargador}}
\label{\detokenize{codigo:module-cargador}}\label{\detokenize{codigo:codigo-de-cargador}}\index{cargador (módulo)@\spxentry{cargador}\spxextra{módulo}}\index{Cargador (clase en cargador)@\spxentry{Cargador}\spxextra{clase en cargador}}

\begin{fulllineitems}
\phantomsection\label{\detokenize{codigo:cargador.Cargador}}\pysigline{\sphinxbfcode{\sphinxupquote{class }}\sphinxcode{\sphinxupquote{cargador.}}\sphinxbfcode{\sphinxupquote{Cargador}}}
\end{fulllineitems}

\index{CargadorTXT (clase en cargador)@\spxentry{CargadorTXT}\spxextra{clase en cargador}}

\begin{fulllineitems}
\phantomsection\label{\detokenize{codigo:cargador.CargadorTXT}}\pysigline{\sphinxbfcode{\sphinxupquote{class }}\sphinxcode{\sphinxupquote{cargador.}}\sphinxbfcode{\sphinxupquote{CargadorTXT}}}
Este metodo se encarga de inicializar la clase, este constructor es heredado de la clase "Cargador"
@param: No tiene parametros
@return: No retorna nada

\end{fulllineitems}



\subsection{Codigo de: \sphinxstylestrong{Conector}}
\label{\detokenize{codigo:module-conector}}\label{\detokenize{codigo:codigo-de-conector}}\index{conector (módulo)@\spxentry{conector}\spxextra{módulo}}\index{Conector (clase en conector)@\spxentry{Conector}\spxextra{clase en conector}}

\begin{fulllineitems}
\phantomsection\label{\detokenize{codigo:conector.Conector}}\pysiglinewithargsret{\sphinxbfcode{\sphinxupquote{class }}\sphinxcode{\sphinxupquote{conector.}}\sphinxbfcode{\sphinxupquote{Conector}}}{\emph{casilla}, \emph{ventana}}{}
Clase genérica de conector para implementar decoradores.
\index{getCasillaFinal() (método de conector.Conector)@\spxentry{getCasillaFinal()}\spxextra{método de conector.Conector}}

\begin{fulllineitems}
\phantomsection\label{\detokenize{codigo:conector.Conector.getCasillaFinal}}\pysiglinewithargsret{\sphinxbfcode{\sphinxupquote{getCasillaFinal}}}{}{}~\begin{quote}\begin{description}
\item[{Devuelve}] \leavevmode
CasillaEstandar -- Referencia a la casilla donde termina el conector.

\end{description}\end{quote}

\end{fulllineitems}

\index{getCasillaInicio() (método de conector.Conector)@\spxentry{getCasillaInicio()}\spxextra{método de conector.Conector}}

\begin{fulllineitems}
\phantomsection\label{\detokenize{codigo:conector.Conector.getCasillaInicio}}\pysiglinewithargsret{\sphinxbfcode{\sphinxupquote{getCasillaInicio}}}{}{}~\begin{quote}\begin{description}
\item[{Devuelve}] \leavevmode
CasillaEstandar -- Referencia a la casilla donde inicia el conector.

\end{description}\end{quote}

\end{fulllineitems}

\index{getVentana() (método de conector.Conector)@\spxentry{getVentana()}\spxextra{método de conector.Conector}}

\begin{fulllineitems}
\phantomsection\label{\detokenize{codigo:conector.Conector.getVentana}}\pysiglinewithargsret{\sphinxbfcode{\sphinxupquote{getVentana}}}{}{}~\begin{quote}\begin{description}
\item[{Devuelve}] \leavevmode
pygame.display -- Referencia a la ventana de pygame.

\end{description}\end{quote}

\end{fulllineitems}

\index{sendConector() (método de conector.Conector)@\spxentry{sendConector()}\spxextra{método de conector.Conector}}

\begin{fulllineitems}
\phantomsection\label{\detokenize{codigo:conector.Conector.sendConector}}\pysiglinewithargsret{\sphinxbfcode{\sphinxupquote{sendConector}}}{}{}
Esta función envía el mismo conector (self) a sus casillas relacionadas de inicio y fin.
:returns:  None -- Procedimiento sin valor de retorno.

\end{fulllineitems}


\end{fulllineitems}

\index{Escalera (clase en conector)@\spxentry{Escalera}\spxextra{clase en conector}}

\begin{fulllineitems}
\phantomsection\label{\detokenize{codigo:conector.Escalera}}\pysiglinewithargsret{\sphinxbfcode{\sphinxupquote{class }}\sphinxcode{\sphinxupquote{conector.}}\sphinxbfcode{\sphinxupquote{Escalera}}}{\emph{casilla}, \emph{ventana}}{}
Esta clase representa un conector entre dos casillas de tablero, y contiene los
atributos y acciones necesarios para su comportamiento en un tablero de juego.
\index{dibujar() (método de conector.Escalera)@\spxentry{dibujar()}\spxextra{método de conector.Escalera}}

\begin{fulllineitems}
\phantomsection\label{\detokenize{codigo:conector.Escalera.dibujar}}\pysiglinewithargsret{\sphinxbfcode{\sphinxupquote{dibujar}}}{}{}
Este método dibuja el conector como una escalera en la ventana de pygame.
:returns:  None -- Procedimiento sin valor de retorno.

\end{fulllineitems}


\end{fulllineitems}

\index{Serpiente (clase en conector)@\spxentry{Serpiente}\spxextra{clase en conector}}

\begin{fulllineitems}
\phantomsection\label{\detokenize{codigo:conector.Serpiente}}\pysiglinewithargsret{\sphinxbfcode{\sphinxupquote{class }}\sphinxcode{\sphinxupquote{conector.}}\sphinxbfcode{\sphinxupquote{Serpiente}}}{\emph{casilla}, \emph{ventana}}{}
Esta clase representa un conector entre dos casillas de tablero, y contiene los
atributos y acciones necesarios para su comportamiento en un tablero de juego.
\index{dibujar() (método de conector.Serpiente)@\spxentry{dibujar()}\spxextra{método de conector.Serpiente}}

\begin{fulllineitems}
\phantomsection\label{\detokenize{codigo:conector.Serpiente.dibujar}}\pysiglinewithargsret{\sphinxbfcode{\sphinxupquote{dibujar}}}{}{}
Este método dibuja el conector como una escalera en la ventana de pygame.
:returns:  None -- Procedimiento sin valor de retorno.

\end{fulllineitems}


\end{fulllineitems}



\subsection{Codigo de: \sphinxstylestrong{Controlador}}
\label{\detokenize{codigo:module-controlador}}\label{\detokenize{codigo:codigo-de-controlador}}\index{controlador (módulo)@\spxentry{controlador}\spxextra{módulo}}\index{empezarCondicionesJuego() (en el módulo controlador)@\spxentry{empezarCondicionesJuego()}\spxextra{en el módulo controlador}}

\begin{fulllineitems}
\phantomsection\label{\detokenize{codigo:controlador.empezarCondicionesJuego}}\pysiglinewithargsret{\sphinxcode{\sphinxupquote{controlador.}}\sphinxbfcode{\sphinxupquote{empezarCondicionesJuego}}}{\emph{controlador}}{}
Llena las condiciones iniciales que ocupa el juego
:param controlador: objeto de controlador para utilizar sus funciones
:type controlador: Controlador object

\end{fulllineitems}

\index{juego() (en el módulo controlador)@\spxentry{juego()}\spxextra{en el módulo controlador}}

\begin{fulllineitems}
\phantomsection\label{\detokenize{codigo:controlador.juego}}\pysiglinewithargsret{\sphinxcode{\sphinxupquote{controlador.}}\sphinxbfcode{\sphinxupquote{juego}}}{\emph{controlador}}{}
Este metodo tiene el ciclo principal del juego, que no termina hasta que alguien gane.
Va a llamar todas las funciones que ocupa el juego para funcionar.
:param controlador: objeto de controlador para utilizar sus funciones
:type controlador: Controlador object

\end{fulllineitems}

\index{run() (en el módulo controlador)@\spxentry{run()}\spxextra{en el módulo controlador}}

\begin{fulllineitems}
\phantomsection\label{\detokenize{codigo:controlador.run}}\pysiglinewithargsret{\sphinxcode{\sphinxupquote{controlador.}}\sphinxbfcode{\sphinxupquote{run}}}{}{}
Metodo que verifica si es una partida cargada o una partida nueva. Tambien despliega el
menu principal. Llama al ciclo principal del juego (metodo juego)

\end{fulllineitems}



\subsection{Codigo de: \sphinxstylestrong{Inventario}}
\label{\detokenize{codigo:module-inventario}}\label{\detokenize{codigo:codigo-de-inventario}}\index{inventario (módulo)@\spxentry{inventario}\spxextra{módulo}}\index{ElementoInventario (clase en inventario)@\spxentry{ElementoInventario}\spxextra{clase en inventario}}

\begin{fulllineitems}
\phantomsection\label{\detokenize{codigo:inventario.ElementoInventario}}\pysigline{\sphinxbfcode{\sphinxupquote{class }}\sphinxcode{\sphinxupquote{inventario.}}\sphinxbfcode{\sphinxupquote{ElementoInventario}}}
\end{fulllineitems}

\index{Ficha (clase en inventario)@\spxentry{Ficha}\spxextra{clase en inventario}}

\begin{fulllineitems}
\phantomsection\label{\detokenize{codigo:inventario.Ficha}}\pysiglinewithargsret{\sphinxbfcode{\sphinxupquote{class }}\sphinxcode{\sphinxupquote{inventario.}}\sphinxbfcode{\sphinxupquote{Ficha}}}{\emph{figura}}{}
\end{fulllineitems}



\subsection{Codigo de: \sphinxstylestrong{Sidebar}}
\label{\detokenize{codigo:module-sidebar}}\label{\detokenize{codigo:codigo-de-sidebar}}\index{sidebar (módulo)@\spxentry{sidebar}\spxextra{módulo}}\index{Sidebar (clase en sidebar)@\spxentry{Sidebar}\spxextra{clase en sidebar}}

\begin{fulllineitems}
\phantomsection\label{\detokenize{codigo:sidebar.Sidebar}}\pysiglinewithargsret{\sphinxbfcode{\sphinxupquote{class }}\sphinxcode{\sphinxupquote{sidebar.}}\sphinxbfcode{\sphinxupquote{Sidebar}}}{\emph{alto\_total}, \emph{ancho\_total}, \emph{ancho\_sidebar}, \emph{mi\_color}, \emph{elementos}}{}
Clase que representa al espacio extra en la ventana de pygame, que no pertenece al tablero.
Contiene informacion valiosa que le comuncia al jugador el turno, quien gana, como moverse...
Ademas, en este espacio se juega con los dados para decidir turnos y permitir al jugador moverse
en el juego.
\index{dibujar() (método de sidebar.Sidebar)@\spxentry{dibujar()}\spxextra{método de sidebar.Sidebar}}

\begin{fulllineitems}
\phantomsection\label{\detokenize{codigo:sidebar.Sidebar.dibujar}}\pysiglinewithargsret{\sphinxbfcode{\sphinxupquote{dibujar}}}{}{}
Metodo que ilustra visualmente el espacio al lado del tablero.

\end{fulllineitems}


\end{fulllineitems}



\subsection{Codigo de: \sphinxstylestrong{Singleton}}
\label{\detokenize{codigo:module-singleton}}\label{\detokenize{codigo:codigo-de-singleton}}\index{singleton (módulo)@\spxentry{singleton}\spxextra{módulo}}\index{Singleton (clase en singleton)@\spxentry{Singleton}\spxextra{clase en singleton}}

\begin{fulllineitems}
\phantomsection\label{\detokenize{codigo:singleton.Singleton}}\pysigline{\sphinxbfcode{\sphinxupquote{class }}\sphinxcode{\sphinxupquote{singleton.}}\sphinxbfcode{\sphinxupquote{Singleton}}}
Esta clase corresponde al cumplimiento del patron de disenyo 
llamado Singletone. Es una clase con metodos de clase, que contiene
informacion de interes para todas las demas clases del programa.
\index{agregar\_jugador() (método de clase de singleton.Singleton)@\spxentry{agregar\_jugador()}\spxextra{método de clase de singleton.Singleton}}

\begin{fulllineitems}
\phantomsection\label{\detokenize{codigo:singleton.Singleton.agregar_jugador}}\pysiglinewithargsret{\sphinxbfcode{\sphinxupquote{classmethod }}\sphinxbfcode{\sphinxupquote{agregar\_jugador}}}{\emph{jugador}}{}
Este metodo reporta la existencia de un jugador en el juego, agregandolo a la lista de jugadores. 
Ademas, le otorga un nombre automaticamente segun sea el orden en el que fue creado.
:param jugador: jugador a agregar.
:type jugador: Jugador

\end{fulllineitems}

\index{establecer\_pos\_jugadores() (método de clase de singleton.Singleton)@\spxentry{establecer\_pos\_jugadores()}\spxextra{método de clase de singleton.Singleton}}

\begin{fulllineitems}
\phantomsection\label{\detokenize{codigo:singleton.Singleton.establecer_pos_jugadores}}\pysiglinewithargsret{\sphinxbfcode{\sphinxupquote{classmethod }}\sphinxbfcode{\sphinxupquote{establecer\_pos\_jugadores}}}{}{}
Este metodo establece las posiciones en las que se van a ubicar los jugadores. Esta pensado
para un minimo de 2 jugadores y un maximo que corresponde a 4. Los ubica en la casilla inicial pero
de manera que se visualicen, anque hayan varios en la misma casilla.

\end{fulllineitems}

\index{getAlto() (método de clase de singleton.Singleton)@\spxentry{getAlto()}\spxextra{método de clase de singleton.Singleton}}

\begin{fulllineitems}
\phantomsection\label{\detokenize{codigo:singleton.Singleton.getAlto}}\pysiglinewithargsret{\sphinxbfcode{\sphinxupquote{classmethod }}\sphinxbfcode{\sphinxupquote{getAlto}}}{}{}
Este metodo devuelve la medida del alto del tablero.
:returns: la medida del alto del tablero

\end{fulllineitems}

\index{getAncho() (método de clase de singleton.Singleton)@\spxentry{getAncho()}\spxextra{método de clase de singleton.Singleton}}

\begin{fulllineitems}
\phantomsection\label{\detokenize{codigo:singleton.Singleton.getAncho}}\pysiglinewithargsret{\sphinxbfcode{\sphinxupquote{classmethod }}\sphinxbfcode{\sphinxupquote{getAncho}}}{}{}
Este metodo devuelve la medida del ancho del tablero.
:returns: la medida del ancho del tablero

\end{fulllineitems}

\index{getCantX() (método de clase de singleton.Singleton)@\spxentry{getCantX()}\spxextra{método de clase de singleton.Singleton}}

\begin{fulllineitems}
\phantomsection\label{\detokenize{codigo:singleton.Singleton.getCantX}}\pysiglinewithargsret{\sphinxbfcode{\sphinxupquote{classmethod }}\sphinxbfcode{\sphinxupquote{getCantX}}}{}{}
Este metodo devuelve la cantidad de casillas horizontales
:returns: la cantidad de casillas en X (horizontalmente)

\end{fulllineitems}

\index{getCantY() (método de clase de singleton.Singleton)@\spxentry{getCantY()}\spxextra{método de clase de singleton.Singleton}}

\begin{fulllineitems}
\phantomsection\label{\detokenize{codigo:singleton.Singleton.getCantY}}\pysiglinewithargsret{\sphinxbfcode{\sphinxupquote{classmethod }}\sphinxbfcode{\sphinxupquote{getCantY}}}{}{}
Este metodo devuelve la cantidad de casillas verticales
:returns: la cantidad de casillas en X (verticalmente)

\end{fulllineitems}

\index{getJugadorGanador() (método de clase de singleton.Singleton)@\spxentry{getJugadorGanador()}\spxextra{método de clase de singleton.Singleton}}

\begin{fulllineitems}
\phantomsection\label{\detokenize{codigo:singleton.Singleton.getJugadorGanador}}\pysiglinewithargsret{\sphinxbfcode{\sphinxupquote{classmethod }}\sphinxbfcode{\sphinxupquote{getJugadorGanador}}}{}{}
Metodo que devuelve al jugador ganador de la partida.

\end{fulllineitems}

\index{getPosJugadores() (método de clase de singleton.Singleton)@\spxentry{getPosJugadores()}\spxextra{método de clase de singleton.Singleton}}

\begin{fulllineitems}
\phantomsection\label{\detokenize{codigo:singleton.Singleton.getPosJugadores}}\pysiglinewithargsret{\sphinxbfcode{\sphinxupquote{classmethod }}\sphinxbfcode{\sphinxupquote{getPosJugadores}}}{\emph{indice}}{}
Este metodo devuelve la posicion en el tablero, en coordenadas, de un jugador en especificion.
:param indice: numero identificador del jugador
:type indice: int
:returns: coordenadas de X y Y del jugador respectivo segun el indice

\end{fulllineitems}

\index{init() (método de clase de singleton.Singleton)@\spxentry{init()}\spxextra{método de clase de singleton.Singleton}}

\begin{fulllineitems}
\phantomsection\label{\detokenize{codigo:singleton.Singleton.init}}\pysiglinewithargsret{\sphinxbfcode{\sphinxupquote{classmethod }}\sphinxbfcode{\sphinxupquote{init}}}{\emph{an}, \emph{al}, \emph{cantX}, \emph{cantY}, \emph{casillaGanadora}, \emph{es\_partida\_carg}, \emph{vent}}{}~\begin{quote}\begin{description}
\item[{Parámetros}] \leavevmode\begin{itemize}
\item {} 
\sphinxstyleliteralstrong{\sphinxupquote{an}} (\sphinxhref{https://docs.python.org/3/library/functions.html\#int}{\sphinxstyleliteralemphasis{\sphinxupquote{int}}}) -- ancho del tablero

\item {} 
\sphinxstyleliteralstrong{\sphinxupquote{al}} (\sphinxhref{https://docs.python.org/3/library/functions.html\#int}{\sphinxstyleliteralemphasis{\sphinxupquote{int}}}) -- largo del tablero.

\item {} 
\sphinxstyleliteralstrong{\sphinxupquote{cantX}} (\sphinxhref{https://docs.python.org/3/library/functions.html\#int}{\sphinxstyleliteralemphasis{\sphinxupquote{int}}}) -- cantidad de casillas horizontalmente (en coordenadas X)

\item {} 
\sphinxstyleliteralstrong{\sphinxupquote{cantX}} -- cantidad de casillas verticalmente (en coordenadas Y)

\item {} 
\sphinxstyleliteralstrong{\sphinxupquote{casillaGanadora}} (\sphinxhref{https://docs.python.org/3/library/functions.html\#int}{\sphinxstyleliteralemphasis{\sphinxupquote{int}}}) -- posicion de casilla ganadora

\item {} 
\sphinxstyleliteralstrong{\sphinxupquote{casillasT}} -- casillas del tablero

\item {} 
\sphinxstyleliteralstrong{\sphinxupquote{vent}} (\sphinxstyleliteralemphasis{\sphinxupquote{pygame.display}}) -- Referencia a la ventana de pygame.

\end{itemize}

\item[{Type}] \leavevmode
vector de objetos Casilla

\end{description}\end{quote}

\end{fulllineitems}

\index{setJugadorGanador() (método de clase de singleton.Singleton)@\spxentry{setJugadorGanador()}\spxextra{método de clase de singleton.Singleton}}

\begin{fulllineitems}
\phantomsection\label{\detokenize{codigo:singleton.Singleton.setJugadorGanador}}\pysiglinewithargsret{\sphinxbfcode{\sphinxupquote{classmethod }}\sphinxbfcode{\sphinxupquote{setJugadorGanador}}}{\emph{jugador}}{}
Metodo que va a informarle al Singleton quien es el jugador ganador. Va a almacenarlo en la variable
jugadorGanador.
:param jugador: Jugador ganador a almacenar en Singleton
:type jugador: Jugador

\end{fulllineitems}

\index{ver\_jugadores() (método de clase de singleton.Singleton)@\spxentry{ver\_jugadores()}\spxextra{método de clase de singleton.Singleton}}

\begin{fulllineitems}
\phantomsection\label{\detokenize{codigo:singleton.Singleton.ver_jugadores}}\pysiglinewithargsret{\sphinxbfcode{\sphinxupquote{classmethod }}\sphinxbfcode{\sphinxupquote{ver\_jugadores}}}{}{}
Este metodo imprime en consola la cantidad de jugadores actuales
:returns: la cantidad de casillas en X (horizontalmente)

\end{fulllineitems}


\end{fulllineitems}



\subsection{Codigo de: \sphinxstylestrong{Tablero}}
\label{\detokenize{codigo:module-tablero}}\label{\detokenize{codigo:codigo-de-tablero}}\index{tablero (módulo)@\spxentry{tablero}\spxextra{módulo}}\index{Tablero (clase en tablero)@\spxentry{Tablero}\spxextra{clase en tablero}}

\begin{fulllineitems}
\phantomsection\label{\detokenize{codigo:tablero.Tablero}}\pysiglinewithargsret{\sphinxbfcode{\sphinxupquote{class }}\sphinxcode{\sphinxupquote{tablero.}}\sphinxbfcode{\sphinxupquote{Tablero}}}{\emph{tamanio}, \emph{ventana}, \emph{ancho\_barra}, \emph{es\_partida\_cargada}}{}
Esta clase representa un tablero de juego, y contiene los
atributos y acciones necesarios para su graficación y comportamiento.
\index{cargarCasillas() (método de tablero.Tablero)@\spxentry{cargarCasillas()}\spxextra{método de tablero.Tablero}}

\begin{fulllineitems}
\phantomsection\label{\detokenize{codigo:tablero.Tablero.cargarCasillas}}\pysiglinewithargsret{\sphinxbfcode{\sphinxupquote{cargarCasillas}}}{}{}
Este metodo va a cear las casillas estandar y las va a guardar dentro del vector de casillas.

\end{fulllineitems}

\index{colorearCasillas() (método de tablero.Tablero)@\spxentry{colorearCasillas()}\spxextra{método de tablero.Tablero}}

\begin{fulllineitems}
\phantomsection\label{\detokenize{codigo:tablero.Tablero.colorearCasillas}}\pysiglinewithargsret{\sphinxbfcode{\sphinxupquote{colorearCasillas}}}{}{}
Esta función colorea las casillas con un algoritmo que impide colores consecutivos. 
:returns:  None -- Procedimiento sin valor de retorno.

\end{fulllineitems}

\index{existeConexion() (método de tablero.Tablero)@\spxentry{existeConexion()}\spxextra{método de tablero.Tablero}}

\begin{fulllineitems}
\phantomsection\label{\detokenize{codigo:tablero.Tablero.existeConexion}}\pysiglinewithargsret{\sphinxbfcode{\sphinxupquote{existeConexion}}}{\emph{indice}}{}
Esta función indica si una casilla ya está conectada con otra.
:param indice: El índice de la casilla en la lista general de casillas.
:type indice: int
:returns:  bool -- Devuelve True si la casilla en ese índice ya está relacionada con un conector.

\end{fulllineitems}

\index{generarCasillas() (método de tablero.Tablero)@\spxentry{generarCasillas()}\spxextra{método de tablero.Tablero}}

\begin{fulllineitems}
\phantomsection\label{\detokenize{codigo:tablero.Tablero.generarCasillas}}\pysiglinewithargsret{\sphinxbfcode{\sphinxupquote{generarCasillas}}}{}{}
Este método genera casillas y sus conectores para un tablero de Serpientes y escaleras de manera aleatoria.
:returns:  None -- Procedimiento sin valor de retorno.

\end{fulllineitems}

\index{generar\_conexiones() (método de tablero.Tablero)@\spxentry{generar\_conexiones()}\spxextra{método de tablero.Tablero}}

\begin{fulllineitems}
\phantomsection\label{\detokenize{codigo:tablero.Tablero.generar_conexiones}}\pysiglinewithargsret{\sphinxbfcode{\sphinxupquote{generar\_conexiones}}}{\emph{casillas}, \emph{cantidad\_serpientes}, \emph{cantidad\_escaleras}}{}
Esta función genera dos listas de tuplas a partir de ciertas posiciones, una para serpientes y otra para escaleras. 
Las tuplas tienen los índices de casillas de inicio y fin de cada serpiente o escalera.
:param casillas: Lista de todas las casillas.
:type casillas: list
:param cantidad\_serpientes: Cantidad deseada de serpientes (n).
:type cantidad\_serpientes: int
:param cantidad\_escaleras: Cantidad deseada de escaleras (m).
:type cantidad\_escaleras: int
:returns:  tuple -- Devuelve como tupla las dos listas con tuplas de posiciones de serpientes y escaleras.

\end{fulllineitems}

\index{generar\_escaleras() (método de tablero.Tablero)@\spxentry{generar\_escaleras()}\spxextra{método de tablero.Tablero}}

\begin{fulllineitems}
\phantomsection\label{\detokenize{codigo:tablero.Tablero.generar_escaleras}}\pysiglinewithargsret{\sphinxbfcode{\sphinxupquote{generar\_escaleras}}}{\emph{casillas}, \emph{cantidad}}{}
Esta función genera tuplas con las posiciones de inicio y fin de n escaleras dadas.
:param casillas: Lista de todas las casillas.
:type casillas: list
:param cantidad: Cantidad deseada de escaleras (n).
:type cantidad: int
:returns:  list -- Devuelve un vector con posiciones de escaleras, pares inicios, impares fines, consecutivos.

\end{fulllineitems}

\index{generar\_serpientes() (método de tablero.Tablero)@\spxentry{generar\_serpientes()}\spxextra{método de tablero.Tablero}}

\begin{fulllineitems}
\phantomsection\label{\detokenize{codigo:tablero.Tablero.generar_serpientes}}\pysiglinewithargsret{\sphinxbfcode{\sphinxupquote{generar\_serpientes}}}{\emph{escaleras}, \emph{casillas}, \emph{cantidad}}{}
Esta función genera tuplas con las posiciones de inicio y fin de n serpientes dadas.
:param escaleras: Lista de escaleras (necesaria para no crear serpientes donde hay escaleras).
:type escaleras: list
:param casillas: Lista de todas las casillas.
:type casillas: list
:param cantidad: Cantidad deseada de serpientes (n).
:type cantidad: int
:returns:  list -- Devuelve un vector con posiciones de serpientes, pares inicios, impares fines, consecutivos.

\end{fulllineitems}

\index{getPar() (método de tablero.Tablero)@\spxentry{getPar()}\spxextra{método de tablero.Tablero}}

\begin{fulllineitems}
\phantomsection\label{\detokenize{codigo:tablero.Tablero.getPar}}\pysiglinewithargsret{\sphinxbfcode{\sphinxupquote{getPar}}}{\emph{indice}, \emph{posiciones}}{}
Esta función devuelve la tupla de casillas correspondiente a una tupla de índices inicio y fin.
:param indice: El índice de la casilla en la lista general de casillas.
:type indice: int
:param posiciones: Lista de tuplas con las posiciones inicial y final de cada serpiente o escalera.
:type posiciones: list
:returns:  tuple -- Devuelve la tupla de dos casillas, la de inicio en pos 0 y la de fin en pos 1.

\end{fulllineitems}

\index{getPosicionesEnTupla() (método de tablero.Tablero)@\spxentry{getPosicionesEnTupla()}\spxextra{método de tablero.Tablero}}

\begin{fulllineitems}
\phantomsection\label{\detokenize{codigo:tablero.Tablero.getPosicionesEnTupla}}\pysiglinewithargsret{\sphinxbfcode{\sphinxupquote{getPosicionesEnTupla}}}{\emph{indice}, \emph{posiciones}}{}
Esta función devuelve la tupla de índices inicio y fin donde se encuentra un índice de casilla.
:param indice: El índice de la casilla en la lista general de casillas.
:type indice: int
:param posiciones: Lista de tuplas con las posiciones inicial y final de cada serpiente o escalera.
:type posiciones: list
:returns:  tuple/None -- Devuelve la tupla si se encontró, si no, devuelve None.

\end{fulllineitems}

\index{obtener\_casillas() (método de tablero.Tablero)@\spxentry{obtener\_casillas()}\spxextra{método de tablero.Tablero}}

\begin{fulllineitems}
\phantomsection\label{\detokenize{codigo:tablero.Tablero.obtener_casillas}}\pysiglinewithargsret{\sphinxbfcode{\sphinxupquote{obtener\_casillas}}}{\emph{x}, \emph{max\_iteracion}}{}
Esta función genera una lista con las coordenadas de cada casilla en el orden lógico de recorrido del tablero. 
:param x: Ancho de la casilla.
:type x: int
:param max\_iteracion: Máxima cantidad de coordenadas por generar.
:type max\_iteracion: int
:returns:  list -- Devuelve una lista con coordenadas que representan centros.

\end{fulllineitems}

\index{obtener\_centro\_casillas() (método de tablero.Tablero)@\spxentry{obtener\_centro\_casillas()}\spxextra{método de tablero.Tablero}}

\begin{fulllineitems}
\phantomsection\label{\detokenize{codigo:tablero.Tablero.obtener_centro_casillas}}\pysiglinewithargsret{\sphinxbfcode{\sphinxupquote{obtener\_centro\_casillas}}}{\emph{ancho\_total}, \emph{alto\_total}}{}
Esta función calcula una cantidad de casillas basado en el tamaño del tablero, y arroja la lista de sus centros. 
:param ancho\_total: Ancho total del tablero.
:type ancho\_total: int
:param alto\_total: Alto total del tablero.
:type alto\_total: int
:returns:  list -- Devuelve una lista con coordenadas que representan centros.

\end{fulllineitems}

\index{obtener\_orden\_coordenadas() (método de tablero.Tablero)@\spxentry{obtener\_orden\_coordenadas()}\spxextra{método de tablero.Tablero}}

\begin{fulllineitems}
\phantomsection\label{\detokenize{codigo:tablero.Tablero.obtener_orden_coordenadas}}\pysiglinewithargsret{\sphinxbfcode{\sphinxupquote{obtener\_orden\_coordenadas}}}{\emph{centro\_casillas\_x}, \emph{centro\_casillas\_y}}{}
Esta función genera una lista con las coordenadas de cada casilla en el orden lógico de recorrido del tablero. 
:param centro\_casillas\_x: Lista de centros de casillas en eje x.
:type centro\_casillas\_x: list
:param centro\_casillas\_y: Lista de centros de casillas en eje y.
:type centro\_casillas\_y: list
:returns:  list -- Devuelve la lista con coordenadas de todas las casillas (x, y) como tuplas.

\end{fulllineitems}


\end{fulllineitems}

\index{VistaTablero (clase en tablero)@\spxentry{VistaTablero}\spxextra{clase en tablero}}

\begin{fulllineitems}
\phantomsection\label{\detokenize{codigo:tablero.VistaTablero}}\pysiglinewithargsret{\sphinxbfcode{\sphinxupquote{class }}\sphinxcode{\sphinxupquote{tablero.}}\sphinxbfcode{\sphinxupquote{VistaTablero}}}{\emph{controlador}, \emph{tamanio}, \emph{ventana}, \emph{ancho\_barra}, \emph{es\_partida\_cargada}}{}
Esta clase representa un tablero de juego, y contiene los
atributos y acciones necesarios para su graficación y comportamiento.
\index{dibujarCasillas() (método de tablero.VistaTablero)@\spxentry{dibujarCasillas()}\spxextra{método de tablero.VistaTablero}}

\begin{fulllineitems}
\phantomsection\label{\detokenize{codigo:tablero.VistaTablero.dibujarCasillas}}\pysiglinewithargsret{\sphinxbfcode{\sphinxupquote{dibujarCasillas}}}{}{}
Esta función dibuja las casillas y sus conectores en la ventana.
:returns:  None -- Procedimiento sin valor de retorno.

\end{fulllineitems}


\end{fulllineitems}



\subsection{Codigo de: \sphinxstylestrong{Vista}}
\label{\detokenize{codigo:module-vista}}\label{\detokenize{codigo:codigo-de-vista}}\index{vista (módulo)@\spxentry{vista}\spxextra{módulo}}

\section{\sphinxstylestrong{Tutorial acerca como jugar Escaleras y Serpientes}}
\label{\detokenize{tutorial:tutorial-acerca-como-jugar-escaleras-y-serpientes}}\label{\detokenize{tutorial::doc}}

\subsection{\sphinxstylestrong{Requerimientos}}
\label{\detokenize{tutorial:requerimientos}}
Para el funcionamiento correcto del programa se requiere tener instalado python 3, al igual que las bibliotecas de pygame y tkinter (en el caso de Windows ya viene instalado con python). Para poder descargar estas bibliotecas se pueden ejecutar los siguientes comandos:

En Linux:

\begin{sphinxVerbatim}[commandchars=\\\{\}]
\PYG{n}{sudo} \PYG{n}{apt} \PYG{n}{install} \PYG{n}{python3}\PYG{o}{\PYGZhy{}}\PYG{n}{setuptools}
\PYG{n}{sudo} \PYG{n}{easy\PYGZus{}install3} \PYG{n}{pip}
\PYG{n}{python3} \PYG{o}{\PYGZhy{}}\PYG{n}{m} \PYG{n}{pip} \PYG{n}{install} \PYG{o}{\PYGZhy{}}\PYG{o}{\PYGZhy{}}\PYG{n}{user} \PYG{n}{pygame}

\PYG{n}{apt}\PYG{o}{\PYGZhy{}}\PYG{n}{get} \PYG{n}{install} \PYG{n}{python}\PYG{o}{\PYGZhy{}}\PYG{n}{tk}
\end{sphinxVerbatim}

En Windows:

\begin{sphinxVerbatim}[commandchars=\\\{\}]
pip install wheel
pip install pygame‑1.9.2a0‑cp34‑none‑win\PYGZus{}AMD64.whl

Para corroborar que se tiene tkinter: python \PYGZhy{}m Tkinter t
\end{sphinxVerbatim}


\subsection{\sphinxstylestrong{Comandos para correr el programa}}
\label{\detokenize{tutorial:comandos-para-correr-el-programa}}
El programa funciona tanto en Windows como en Linux, y debe de estarse ubicado en el directorio principal de Escaleras\_y\_Serpientes para poder correr el juego. El comando para correr el programa en Linux es:

\begin{sphinxVerbatim}[commandchars=\\\{\}]
\PYG{n}{python3} \PYG{n}{source}\PYG{o}{/}\PYG{n}{controlador}\PYG{o}{.}\PYG{n}{py}
\end{sphinxVerbatim}

Si se desea correr en Windows, se usa con el comando:

\begin{sphinxVerbatim}[commandchars=\\\{\}]
\PYG{n}{python} \PYG{n}{source}\PYG{o}{/}\PYG{n}{controlador}\PYG{o}{.}\PYG{n}{py}
\end{sphinxVerbatim}


\subsection{\sphinxstylestrong{Instrucciones para empezar a jugar}}
\label{\detokenize{tutorial:instrucciones-para-empezar-a-jugar}}
Cuando se inicia el programa, se va abrir un menú con las siguientes opciones:

\begin{sphinxVerbatim}[commandchars=\\\{\}]
\PYG{n}{Cargar} \PYG{n}{partida}
\PYG{n}{Empezar} \PYG{n}{partida}
\end{sphinxVerbatim}

Si se escoge \sphinxstylestrong{cargar partida}
\begin{itemize}
\item {} 
Se tiene que escribir el nombre del archivo de texto de la partida guardada. No se aceptan archivos de texto vacios o inexistentes. Si el archivo tiene menos líneas de texto del tamaño mínimo, que es 7, se considera corrupto y el programa tira excepción para evitar cear una partida con datos incorrectos.

\end{itemize}

Los archivos de texto que funcionan como guardados, se almacenan en la carpeta \sphinxstylestrong{Guardados} y tienen este formato:

\begin{sphinxVerbatim}[commandchars=\\\{\}]
\PYG{n}{Cantidad} \PYG{n}{de} \PYG{n}{jugadores}
\PYG{n}{Nombre} \PYG{k}{del} \PYG{n}{jugador} \PYG{o}{/} \PYG{n}{posicion} \PYG{n}{en} \PYG{n}{el} \PYG{n}{tablero} \PYG{o}{/} \PYG{n}{color} \PYG{k}{del} \PYG{n}{jugador} \PYG{o}{/} \PYG{n}{posicion} \PYG{n}{dentro} \PYG{n}{de} \PYG{n}{una} \PYG{n}{casilla} \PYG{n}{en} \PYG{n}{X} \PYG{o}{/} \PYG{n}{posicion} \PYG{n}{dentro} \PYG{n}{de} \PYG{n}{una} \PYG{n}{casilla} \PYG{n}{en} \PYG{n}{X}
\PYG{o}{\PYGZhy{}}\PYG{o}{\PYGZhy{}}\PYG{o}{\PYGZhy{}}\PYG{n}{Hasta} \PYG{n}{la} \PYG{n}{cantidad} \PYG{n}{de} \PYG{n}{jugadores}\PYG{o}{.}\PYG{o}{.}\PYG{o}{.}
\PYG{n}{Cantidad} \PYG{n}{de} \PYG{n}{serpientes}
\PYG{n}{Posicion} \PYG{n}{en} \PYG{n}{el} \PYG{n}{tablero} \PYG{n}{de} \PYG{n}{la} \PYG{n}{casilla} \PYG{n}{de} \PYG{n}{inicio} \PYG{n}{de} \PYG{n}{serpiente} \PYG{l+m+mi}{1} \PYG{o}{/} \PYG{n}{posicion} \PYG{n}{en} \PYG{n}{el} \PYG{n}{tablero} \PYG{n}{de} \PYG{n}{la} \PYG{n}{casilla} \PYG{n}{final} \PYG{n}{de} \PYG{n}{serpiente} \PYG{l+m+mi}{1}
\PYG{o}{.}\PYG{o}{.}\PYG{o}{.} \PYG{n}{Hasta} \PYG{n}{la} \PYG{n}{cantidada} \PYG{n}{de} \PYG{n}{serpientes} \PYG{o}{.}\PYG{o}{.}\PYG{o}{.}
\PYG{n}{Cantidad} \PYG{n}{de} \PYG{n}{escaleras}
\PYG{n}{Posicion} \PYG{n}{en} \PYG{n}{el} \PYG{n}{tablero} \PYG{n}{de} \PYG{n}{la} \PYG{n}{casilla} \PYG{n}{de} \PYG{n}{inicio} \PYG{n}{de} \PYG{n}{escalera} \PYG{l+m+mi}{1} \PYG{o}{/} \PYG{n}{posicion} \PYG{n}{en} \PYG{n}{el} \PYG{n}{tablero} \PYG{n}{de} \PYG{n}{la} \PYG{n}{casilla} \PYG{n}{final} \PYG{n}{de} \PYG{n}{escalera} \PYG{l+m+mi}{1}
\PYG{o}{.}\PYG{o}{.}\PYG{o}{.} \PYG{n}{Hasta} \PYG{n}{la} \PYG{n}{cantidada} \PYG{n}{de} \PYG{n}{escaleras} \PYG{o}{.}\PYG{o}{.}\PYG{o}{.}
\end{sphinxVerbatim}

Si se escoge \sphinxstylestrong{empezar partida}
\begin{itemize}
\item {} 
El programa va a solicitarle al usuario que se digite el numero de jugadores que van a estar en la partida y los colores para sus fichas. Si estos datos no se ingresan, el juego no va a iniciar. Solo se pueden entre 2 a 4 jugadores en la partida y no se pueden repetir colores. Cualquier ingreso inválido o que no cumpla los requisitos recíen mencionados, ocasionará que el programa le solicite la información denuevo e incluso mencionará su error.

\end{itemize}


\subsection{\sphinxstylestrong{Reglas del juego}}
\label{\detokenize{tutorial:reglas-del-juego}}
Las reglas para comprender el juego antes de iniciarlo, son:

Número de jugadores: dos o más, pero nunca deben jugar más de cuatro jugadores a la vez.
Las lineas rojas corresponden a serpientes y las verdes a escaleras.
Se necesita: un tablero del juego, una ficha de diferente color para cada jugador y un dado.

Reglas dentro del juego:
1- Para comenzar, todos los jugadores tenéis que tirar una vez el dado. El que saque más puntos será el primero en empezar la partida; continuará el que esté a su izquierda y así sucesivamente.
2- El juego consiste en mover la ficha desde la casilla 1 hasta la 100; debes avanzarla tantos puestos como puntos te salgan al tirar el dado.
3- Pero ¡cuidado!, si caes en la cabeza de una serpiente debes bajar hasta su cola. LAS SERPIENTES SOLO BAJAN, NUNCA SUBEN.
4- También puedes tener suerte. Si caes en una casilla que señale el principio de una escalera sube hasta el final de la escalera. LAS ESCALERAS SOLO SUBEN, NUNCA BAJAN.
5- Gana el primero que llegue a la última casilla.


\section{\sphinxstylestrong{Informacion acerca del proyecto: Escaleras y Serpientes}}
\label{\detokenize{proyecto:informacion-acerca-del-proyecto-escaleras-y-serpientes}}\label{\detokenize{proyecto::doc}}

\subsection{\sphinxstylestrong{Autores}}
\label{\detokenize{proyecto:autores}}\begin{itemize}
\item {} 
Daniel Salazar

\item {} 
Ricardo Franco

\item {} 
Gustavo Jimenez

\item {} 
Ignacio Arroyo

\item {} 
Josue Amador

\end{itemize}


\subsection{\sphinxstylestrong{Descripcion del proyecto}}
\label{\detokenize{proyecto:descripcion-del-proyecto}}
Trabajo final del juego Escaleras y Serpientes implementado como MARDA.


\chapter{\sphinxstylestrong{Indices}}
\label{\detokenize{Escaleras_y_Serpientes:indices}}\begin{itemize}
\item {} 
\DUrole{xref,std,std-ref}{genindex}

\item {} 
\DUrole{xref,std,std-ref}{modindex}

\item {} 
\DUrole{xref,std,std-ref}{search}

\end{itemize}


\renewcommand{\indexname}{Índice de Módulos Python}
\begin{sphinxtheindex}
\let\bigletter\sphinxstyleindexlettergroup
\bigletter{a}
\item\relax\sphinxstyleindexentry{almacenador}\sphinxstyleindexpageref{codigo:\detokenize{module-almacenador}}
\item\relax\sphinxstyleindexentry{arbitro}\sphinxstyleindexpageref{codigo:\detokenize{module-arbitro}}
\indexspace
\bigletter{b}
\item\relax\sphinxstyleindexentry{button}\sphinxstyleindexpageref{codigo:\detokenize{module-button}}
\indexspace
\bigletter{c}
\item\relax\sphinxstyleindexentry{cargador}\sphinxstyleindexpageref{codigo:\detokenize{module-cargador}}
\item\relax\sphinxstyleindexentry{conector}\sphinxstyleindexpageref{codigo:\detokenize{module-conector}}
\item\relax\sphinxstyleindexentry{controlador}\sphinxstyleindexpageref{codigo:\detokenize{module-controlador}}
\indexspace
\bigletter{d}
\item\relax\sphinxstyleindexentry{dado}\sphinxstyleindexpageref{codigo:\detokenize{module-dado}}
\indexspace
\bigletter{i}
\item\relax\sphinxstyleindexentry{inventario}\sphinxstyleindexpageref{codigo:\detokenize{module-inventario}}
\indexspace
\bigletter{j}
\item\relax\sphinxstyleindexentry{jugador}\sphinxstyleindexpageref{codigo:\detokenize{module-jugador}}
\indexspace
\bigletter{s}
\item\relax\sphinxstyleindexentry{sidebar}\sphinxstyleindexpageref{codigo:\detokenize{module-sidebar}}
\item\relax\sphinxstyleindexentry{singleton}\sphinxstyleindexpageref{codigo:\detokenize{module-singleton}}
\indexspace
\bigletter{t}
\item\relax\sphinxstyleindexentry{tablero}\sphinxstyleindexpageref{codigo:\detokenize{module-tablero}}
\indexspace
\bigletter{v}
\item\relax\sphinxstyleindexentry{vista}\sphinxstyleindexpageref{codigo:\detokenize{module-vista}}
\end{sphinxtheindex}

\renewcommand{\indexname}{Índice}
\printindex
\end{document}